\begin{otherlanguage}{english}
Scientific knowledge about the medicinal plants of Burkina Faso is scattered across articles, theses and reports that use varying terminology. As a result, it is difficult to find a piece of information and trace it back to its source. This thesis proposes intelligent and traceable access to this knowledge.

After comparing five ontologies, we selected \textit{ontoMEDTRAD} and adapted its conceptual model to the Burkinabè context to build \textit{OntoPlantMedBF}. We populated it from ten of the 35 documents in the corpus, using extraction assisted by a language model. The ontology contains 92 classes and 1,546 individuals, including 141 plants and 277 therapeutic uses, all linked to a source. Imported into Neo4j, it feeds a Retrieval-Augmented Generation (RAG) system: fixed queries select the facts from the graph, then Llama 3 writes an answer that cites its sources. A web application built with Streamlit provides access to the whole system.

On a set of 50 questions, the system reaches an F1-score of 0.923 and produces 42 compliant answers. It declines to answer the 5 out-of-corpus questions. These measurements concern how well the content of the graph is restored, not the scientific accuracy of the facts, as the extractions were not reviewed by an expert. This prototype is a consultation tool: it does not make diagnoses or issue prescriptions.

\vspace{1em}
\noindent\textbf{Keywords:} medicinal plants, Burkina Faso, ontology, knowledge graph, Neo4j, Retrieval-Augmented Generation, language model, traceability.
\end{otherlanguage}
